% !TeX root = ../main.tex
 
\chapter{Template}

\section{Reference Examples}

\begin{Definition}{Referenceable Definition}{referenceable-definition}
	A labeled definition uses the third argument as the label key.
\end{Definition}

\begin{Lemma}{Referenceable Lemma}{referenceable-lemma}
	A labeled lemma works the same way.
\end{Lemma}

\begin{theorem}{Referenceable Theorem}{referenceable-theorem}
	A labeled theorem also works the same way.
\end{theorem}

\Cref{def:referenceable-definition}, \Cref{th:referenceable-lemma}, and
\Cref{th:referenceable-theorem} are clickable references.

\begin{latexcodeblock}[title={Reference Syntax}]
\begin{Definition}{Referenceable Definition}{referenceable-definition}
Body
\end{Definition}

\begin{Lemma}{Referenceable Lemma}{referenceable-lemma}
Body
\end{Lemma}

\begin{theorem}{Referenceable Theorem}{referenceable-theorem}
Body
\end{theorem}

\Cref{def:referenceable-definition}
\Cref{th:referenceable-lemma}
\Cref{th:referenceable-theorem}
\end{latexcodeblock}

\cor{}{By the result of the proof, we can then show...}
\mlemma{}{Suppose $\vec{v_1}, \dots, \vec{v_n} \in \RR[n]$ is subspace of $\RR^n$.}
\mprop{}{$1 + 1 = 2$.}

